%% =====================================================================
%% Final Report -- QAD-EDCA: A Lightweight Queue-Aware Dynamic EDCA
%% Controller for Starvation Avoidance in IEEE 802.11 WLANs
%% IEEE conference format. Compile on Overleaf (pdfLaTeX).
%%
%% NOTE: upload these 3 figure PDFs into the Overleaf project, in a "figures/"
%% folder OR next to this .tex (the \graphicspath below looks in both):
%%   fig_dynamic_hi_delay.pdf, fig_dynamic_be_throughput.pdf, positioning_tradeoff.pdf
%% Do NOT rely on ../analysis/figures/ on Overleaf -- parent-of-project paths are
%% not resolved there. Each figure is also mirrored by an in-text table, so the
%% paper stays complete even if you comment out the \includegraphics lines.
%% =====================================================================
\documentclass[conference]{IEEEtran}
\IEEEoverridecommandlockouts

\usepackage{cite}
\usepackage{amsmath,amssymb,amsfonts}
\usepackage{algorithm}
\usepackage{algpseudocode}
\usepackage{graphicx}
\usepackage{booktabs}
\usepackage{array}
\usepackage{textcomp}
\usepackage{xcolor}
\usepackage{url}
\usepackage[hidelinks]{hyperref}

% Relax float placement: this paper is table-heavy, so let more floats sit per
% page (avoids tables piling up and spilling past the references).
\setcounter{topnumber}{4}
\setcounter{totalnumber}{6}
\renewcommand{\topfraction}{0.95}
\renewcommand{\bottomfraction}{0.85}
\renewcommand{\textfraction}{0.07}
\renewcommand{\floatpagefraction}{0.8}

\graphicspath{{figures/}{./}}

% Compact helper for access-category names (text mode).
\newcommand{\AC}[1]{AC\_#1}

\begin{document}

\title{QAD-EDCA: A Lightweight Queue-Aware Dynamic EDCA Controller for
Starvation Avoidance in IEEE~802.11 WLANs}

\author{
\IEEEauthorblockN{Cheng-Yu Wang (M143140017) \quad Hsien-Yi Lin (M143040081)}
\IEEEauthorblockA{CSE625 --- Wireless Communication Network:
Construction and Performance Simulation\\
National Sun Yat-sen University, Spring 2026}
}

\maketitle

\begin{abstract}
The Enhanced Distributed Channel Access (EDCA) mechanism of IEEE~802.11
provides Quality-of-Service (QoS) differentiation through fixed per-Access-Category
(AC) parameters. Under saturated load these static parameters cause severe
\emph{starvation} of the low-priority Best-Effort (\AC{BE}) and
Background (\AC{BK}) categories. We present QAD-EDCA (Queue-Aware Dynamic
EDCA), a lightweight, rule-based, $O(1)$ controller that runs at the Access
Point, detects starvation from queue occupancy and packet-loss rate, and
adjusts AIFSN, CWmin, and TXOP across all four ACs in closed loop, with
exponential recovery, a QoS safety valve, and Schmitt-trigger hysteresis.
In OMNeT++/INET simulations QAD-EDCA raises starved \AC{BE} throughput by
$2.15\times$ over standard EDCA (2.53$\rightarrow$5.44~Mbps) with no loss of
\AC{VO}/\AC{VI} QoS and \emph{zero} manual tuning or training. Equally important, we report an honest negative result: a
queue-feedback controller \emph{cannot} exceed a well-chosen feedforward
static optimum in saturated steady state. We establish this against a
well-tuned static optimum that beats QAD-EDCA by $\sim$6\% consistently across
load mixes, and we explain the
mechanism: relieving starvation drains the queue that triggers the
controller, so the boost decays (the trigger rate falls from 64\% to
$\sim$15\%). We also document an instructive engineering pitfall---a static
TXOP misconfiguration that produced ``fake dynamics''---and how vector-trace
verification exposed it. QAD-EDCA is therefore positioned as the best option
among methods that need neither prior knowledge of the load nor training,
trailing only a foreknowledge-dependent static optimum and training-dependent
deep reinforcement learning.
\end{abstract}

\begin{IEEEkeywords}
IEEE 802.11, EDCA, QoS, starvation avoidance, contention window, adaptive
MAC, feedback control, WLAN.
\end{IEEEkeywords}

% =====================================================================
\section{Introduction}
% =====================================================================
The Enhanced Distributed Channel Access (EDCA) mechanism in IEEE~802.11
provides Quality-of-Service (QoS) differentiation by classifying traffic into
four Access Categories (ACs)---Voice (\AC{VO}), Video (\AC{VI}), Best Effort
(\AC{BE}), and Background (\AC{BK})---each assigned distinct channel-access
parameters: the Arbitration Inter-Frame Space Number (AIFSN), the contention
window bounds (CWmin/CWmax), and the Transmission Opportunity limit (TXOP)
\cite{ieee80211}. These parameters are negotiated at association and held
constant for the session.

Their static nature produces a well-documented pathology: under
moderate-to-heavy load, high-priority ACs monopolize the channel and the
low-priority \AC{BE}/\AC{BK} categories \emph{starve}, exhibiting near-zero
throughput and excessive loss \cite{ugwu,mammeri}. The protocol contains no
feedback path to detect or relieve this condition.

\textbf{System and problem.} We consider a single infrastructure Basic
Service Set (BSS): one Access Point (AP) and $N$ associated stations carrying
a realistic mix of VO/VI/BE/BK traffic. The practical problem is to relieve
low-priority starvation \emph{at runtime}, without retraining or
operator-specific manual tuning, while never violating the delay budget of
the real-time \AC{VO}/\AC{VI} flows.

\textbf{Approach.} We propose QAD-EDCA, a queue-aware, rule-based controller
that runs only at the AP (which already disseminates the EDCA Parameter Set
in beacons, so stations need no modification). Once per beacon interval it
reads the four queue occupancies and loss rates, evaluates a formal
starvation predicate, and applies bounded, deterministic adjustments to
AIFSN/CWmin/TXOP, with exponential recovery toward defaults and a QoS safety
valve. Its complexity is $O(1)$ per cycle and it acts from the first
monitoring interval---no training.

\textbf{Contributions.}
\begin{enumerate}
\item \emph{A formal, implementable starvation criterion.} Prior work
describes starvation only qualitatively; we give a quantifiable dual-indicator
predicate (queue occupancy OR loss rate) with a justification for the
OR-logic, directly usable as a detector.
\item \emph{A complete $O(1)$ closed-loop controller.} QAD-EDCA adjusts three
parameters governing three distinct phases of channel access, adds
bidirectional CWmin shaping, exponential recovery, a QoS safety valve, and
Schmitt-trigger hysteresis against oscillation.
\item \emph{An honest, rigorous negative result.} We show that a
queue-feedback controller cannot beat a well-chosen feedforward static optimum
in saturated steady state. We confirm this against a carefully hand-tuned
static optimum---the gap is stable across all load mixes---and give the
mechanism: relieving starvation removes the queue signal that triggers the
controller.
We additionally document a methodological pitfall---a static-TXOP
``fake dynamics'' confound---and how vector tracing caught it.
\item \emph{A clear positioning.} Among methods requiring neither foreknowledge
of the load nor training, QAD-EDCA is the strongest: $2.15\times$ \AC{BE} gain
over standard EDCA with no QoS loss, trailing only a foreknowledge-dependent
static optimum and training-dependent DRL.
\end{enumerate}

% =====================================================================
\section{Related Work}\label{sec:related}
% =====================================================================
\textbf{Static and tuned EDCA.} The default EDCA parameter set
\cite{ieee80211} trades fairness for differentiation; Ugwu \emph{et al.}
\cite{ugwu} confirmed via simulation that AIFSN exerts the strongest influence
on the resulting disparity and that low-priority loss rises sharply under
saturation. A hand-tuned static configuration (shorter CWmin/AIFSN for
\AC{BE}/\AC{BK}) can relieve much of the starvation, but only for the load it
was tuned against.

\textbf{Dynamic and multichannel EDCA.} Mammeri \emph{et al.} \cite{mammeri}
proposed a starvation-avoidance dynamic multichannel access scheme for
802.11ac, but it relies on channel bonding and does not apply to a
single-channel BSS. Many adaptive-EDCA proposals adjust only the contention
window and lack a formal starvation criterion.

\textbf{Learning-based EDCA.} Recent deep reinforcement learning (DRL) work is
strong on raw throughput: Zuo \emph{et al.}'s PDCF-DRL \cite{zuo} keeps
normalized throughput above 76\% and collision rate below 18\% across 20--120
stations, far above standard EDCA. Du
\emph{et al.} \cite{du} combine federated learning with DDPG; Li \emph{et al.}
\cite{li} optimize QoS at the application layer. All require training (and,
for DRL, possibly sub-EDCA performance before convergence) or bypass the MAC
layer.

\textbf{Wi-Fi 6/7 context.} EDCA remains foundational in 802.11ax \cite{ax}
and 802.11be \cite{yi}, where much effort targets OBSS/inter-BSS interference
and multi-link operation \cite{tuan}. Lightweight MAC-layer mechanisms for
intra-BSS AC starvation remain comparatively underexplored---the gap this work
addresses.

% =====================================================================
\section{System Model}\label{sec:model}
% =====================================================================
\textbf{Deployment.} We model a single infrastructure BSS in an Indoor-Hotspot
(InH) office of $120\,\text{m}\times50\,\text{m}$ (3GPP TR~38.901
Table~7.2-2): one AP plus $N\in\{5,10,15,20\}$ stations, all using the OFDM
PHY at 54~Mbps ($a\text{SlotTime}=9\,\mu s$, $a\text{SIFSTime}=16\,\mu s$).
Each station hosts one AC; traffic flows uplink from the stations through the
AP to a sink. The AP runs the EDCA Parameter Set advertisement, so it is the
natural and standards-compatible control point---stations are not modified.

\textbf{Traffic model} (standards-aligned; Table~\ref{tab:traffic}). \AC{VO}
follows G.711 over RTP (160~B every 20~ms, 64~kbps; 3GPP 5QI=1).
\AC{VI} models conversational HD video (1280~B every 10~ms,
$\approx$1~Mbps; 5QI=2, PDB~150~ms). \AC{BE} is a saturated bulk UDP source
(1500~B, IEEE~802.3 MTU; 5QI=9). \AC{BK} is background sync (512~B every
100~ms; 5QI=8).

\begin{table}[t]
\centering
\caption{Standards-aligned traffic model.}
\label{tab:traffic}
\begin{tabular}{@{}llrrl@{}}
\toprule
AC & Application & Size & Interval & Rate\\
\midrule
\AC{VO} & VoIP (G.711)        & 160\,B  & 20\,ms  & 64\,kbps\\
\AC{VI} & Conv. HD video      & 1280\,B & 10\,ms  & $\approx$1\,Mbps\\
\AC{BE} & Saturated bulk UDP  & 1500\,B & sat.    & saturated\\
\AC{BK} & Background sync      & 512\,B  & 100\,ms & 40.96\,kbps\\
\bottomrule
\end{tabular}
\end{table}

\textbf{Standard EDCA parameters.} For the OFDM PHY ($a\text{CWmin}=15$,
$a\text{CWmax}=1023$) the default per-AC parameters used as our baseline are
shown in Table~\ref{tab:edca}. We use the INET/IEEE single-MSDU TXOP defaults
($\text{TXOP}=0$ for \AC{BE}/\AC{BK}); the relevance of this choice is
revisited in Sec.~\ref{sec:analysis} (the ``fake dynamics'' lesson).

\begin{table}[t]
\centering
\caption{Default EDCA parameter set (OFDM PHY) used as baseline.}
\label{tab:edca}
\begin{tabular}{@{}lrrrr@{}}
\toprule
AC & CWmin & CWmax & AIFSN & TXOP\\
\midrule
\AC{BK} & 15 & 1023 & 7 & 0\\
\AC{BE} & 15 & 1023 & 3 & 0\\
\AC{VI} & 7  & 15   & 2 & 3.008\,ms\\
\AC{VO} & 3  & 7    & 2 & 1.504\,ms\\
\bottomrule
\end{tabular}
\end{table}

\textbf{Controller observables and scope.} The controller observes, per AC at
the AP: the pending-queue occupancy ratio $Q_i/Q_{cap}$ and the per-interval
loss rate $P_{loss,i}$; it actuates AIFSN, CWmin, and TXOP. By design it
adjusts only the AP-side EDCA (mirroring beacon-disseminated parameters);
station-side parameters are left to the standard mechanism. This scope is a
deliberate simplification and a source of one limitation analyzed later.

% =====================================================================
\section{Problem Formulation}\label{sec:problem}
% =====================================================================
\textbf{Starvation predicate.} Let $AC_i$, $i\in\{\text{BE},\text{BK}\}$, have
queue length $Q_i$, queue capacity $Q_{cap}$, and interval loss rate
$P_{loss,i}$. We define
\begin{equation}
\text{Starv}(AC_i)\iff\Big(\tfrac{Q_i}{Q_{cap}}>Q_{th}\Big)\ \lor\
\big(P_{loss,i}>P_{th}\big),
\label{eq:starv}
\end{equation}
with defaults $Q_{th}=0.8$, $P_{th}=0.1$. The two indicators are
complementary: queue occupancy captures sustained access deprivation (packets
accumulate faster than they drain), while loss rate captures starvation that
manifests as drops---tail-drop at a full queue, or discard after exceeding the
retry limit due to repeated collisions with high-priority traffic. Either
alone misses a failure mode; the OR combination is robust.

\textbf{Why static EDCA is suboptimal.} Default parameters
(Table~\ref{tab:edca}) are fixed at association. The per-slot transmit
probability under a CWmin of $W$ is approximately
\begin{equation}
\tau \approx \frac{2}{W+1},
\label{eq:tau}
\end{equation}
so \AC{VO} ($W{=}3$) transmits with $\tau\approx0.5$ versus \AC{BE}
($W{=}15$) at $\tau\approx0.125$. Under saturation this gap drives \AC{BE}/\AC{BK}
to near-zero throughput. There is no feedback to correct it.

\textbf{Difficulty.} The problem is harder than ``lower the low-priority
CWmin'' for three coupled reasons. (i) \emph{Coupling with QoS:} any boost to
\AC{BE}/\AC{BK} steals airtime from delay-bounded \AC{VO}/\AC{VI}, so the
controller must respect a hard delay budget. (ii) \emph{Oscillation:} a naive
trigger that fully reverts on relief re-creates starvation immediately,
producing ping-pong. (iii) \emph{Observability vs.\ actuation---a structural
trap:} the very act of relieving starvation \emph{removes the signal that
detected it}. A queue-occupancy trigger that successfully drains the queue
then reads ``not starving'' and backs off, so a pure feedback controller is
pulled toward an operating point strictly interior to what a feedforward
optimum (which never backs off) can hold. Sec.~\ref{sec:analysis} shows this
trap is the binding constraint, not an implementation detail.

% =====================================================================
\section{Solution: QAD-EDCA}\label{sec:solution}
% =====================================================================
\subsection{Architecture}
QAD-EDCA is a centralized monitor-controller at the AP. Every
$T_{mon}=100$~ms (aligned with the beacon interval, so no extra overhead) it
reads four queue lengths and four loss rates, evaluates \eqref{eq:starv}, and
applies deterministic adjustments---an $O(1)$ rule engine, in contrast to
$O(\text{DRL inference})$ learned policies that also require offline training.

\subsection{Three-stage adjustment}
On a positive starvation verdict, three complementary actions are applied
\emph{simultaneously}, because they act on three different phases of channel
access: AIFSN governs \emph{when} contention starts, CWmin \emph{how long}
backoff lasts, and TXOP \emph{how long} the channel is held after winning.

\emph{Stage 1 --- lower AIFSN of low-priority ACs:}
\begin{equation}
\text{AIFSN}_i=\max(\text{AIFSN}_i-\Delta_{AIFS},\,2),\quad i\in\{\text{BE},\text{BK}\}.
\end{equation}
The floor of 2 (the \AC{VO}/\AC{VI} default) ensures low-priority ACs never
become \emph{more} aggressive than high-priority ones. AIFSN is the strongest
single lever \cite{ugwu}.

\emph{Stage 2 --- shape CWmin (bidirectional):} raise high-priority CWmin to
moderate dominance and lower low-priority CWmin to accelerate relief,
\begin{align}
\text{CWmin}_i &=\min(\text{CWmin}_i\cdot\alpha_{CW},\,\text{CWmin}_{BE}^{def}),
& i&\in\{\text{VO},\text{VI}\},\\
\text{CWmin}_j &=\max(\text{CWmin}_j/\alpha_{CW},\,\text{CWmin}_{VI}^{def}),
& j&\in\{\text{BE},\text{BK}\}.
\end{align}
The ceiling $\text{CWmin}_{BE}^{def}{=}15$ and floor
$\text{CWmin}_{VI}^{def}{=}7$ preserve the priority ordering by construction.

\emph{Stage 3 --- reduce TXOP of high-priority ACs:}
\begin{equation}
\text{TXOP}_i=\max(\text{TXOP}_i\cdot\beta_{TXOP},\,\text{TXOP}_{min}),\quad i\in\{\text{VO},\text{VI}\},
\end{equation}
forcing earlier channel release so low-priority traffic gets access.

\subsection{Recovery, safety valve, hysteresis}
\emph{Exponential recovery.} When \eqref{eq:starv} is false, every parameter
relaxes toward its default,
\begin{equation}
P(t{+}1)=P(t)+\gamma\,(P_{def}-P(t)),\quad 0<\gamma<1,
\end{equation}
with $\gamma=0.3$ (\textasciitilde83\% recovery within 5 cycles). Gradual
recovery prevents the ping-pong of instantaneous reversion.

\emph{QoS safety valve.} After adjusting, if a high-priority delay bound is
exceeded ($d_{VO}>150$~ms or $d_{VI}>150$~ms), parameters are pulled halfway
back to default,
$P\leftarrow(P_{adjusted}+P_{def})/2$, so QoS preservation strictly dominates
starvation relief.

\emph{Schmitt-trigger hysteresis.} To avoid chattering, detection switches ON
at the thresholds in \eqref{eq:starv} but OFF only below a lower release band
($0.4\times$ the thresholds), so genuine relief stops throttling while
transient dips do not.

\subsection{Algorithm and parameters}
Algorithm~\ref{alg:qad} summarizes one control cycle; floats accumulate in
double precision to avoid integer-truncation deadlock during slow recovery.
Table~\ref{tab:params} lists the parameters. The five algorithm constants are
design choices, not standardized values; their robustness is established by
the sensitivity sweep in Sec.~\ref{sec:analysis}.

\begin{algorithm}[t]
\caption{QAD-EDCA control cycle (run at the AP every $T_{mon}$)}
\label{alg:qad}
\begin{algorithmic}[1]
\State \textbf{for} $i\in\{\text{BE},\text{BK}\}$: read $Q_i/Q_{cap}$, $P_{loss,i}$
\If{$\text{Starv}(\text{BE})\ \lor\ \text{Starv}(\text{BK})$}
  \State $\text{AIFSN}_{\text{BE},\text{BK}}\gets\max(\text{AIFSN}-\Delta_{AIFS},2)$
  \State $\text{CWmin}_{\text{VO},\text{VI}}\gets\min(\text{CWmin}\cdot\alpha_{CW},15)$
  \State $\text{CWmin}_{\text{BE},\text{BK}}\gets\max(\text{CWmin}/\alpha_{CW},7)$
  \State $\text{TXOP}_{\text{VO},\text{VI}}\gets\max(\text{TXOP}\cdot\beta_{TXOP},\text{TXOP}_{min})$
\Else
  \ForAll{parameters $P$}
    \State $P\gets P+\gamma\,(P_{def}-P)$ \Comment{exponential recovery}
  \EndFor
\EndIf
\State Disseminate updated EDCA set via beacon
\If{$d_{VO}>150$\,ms $\lor\ d_{VI}>150$\,ms}
  \State $P\gets (P+P_{def})/2$ \Comment{QoS safety valve}
\EndIf
\end{algorithmic}
\end{algorithm}

\begin{table}[t]
\centering
\caption{QAD-EDCA parameters.}
\label{tab:params}
\begin{tabular}{@{}lll@{}}
\toprule
Parameter & Symbol & Default\\
\midrule
Monitoring interval & $T_{mon}$ & 100\,ms\\
Queue threshold & $Q_{th}$ & 0.80\\
Loss threshold & $P_{th}$ & 0.10\\
AIFS step & $\Delta_{AIFS}$ & 1 (floor 2)\\
CW scale & $\alpha_{CW}$ & 2.0\\
TXOP scale & $\beta_{TXOP}$ & 0.75\\
Recovery factor & $\gamma$ & 0.30\\
Hysteresis release & --- & $0.4\times$ thresholds\\
\bottomrule
\end{tabular}
\end{table}

\subsection{Fairness metric}
Cross-AC fairness is quantified by Jain's index \cite{jain}
$J(\mathbf{x})=\big(\sum_i x_i\big)^2/\big(n\sum_i x_i^2\big)\in[1/n,1]$
over normalized per-AC throughputs $x_i$.

% =====================================================================
\section{Simulation}\label{sec:sim}
% =====================================================================
\textbf{Platform.} OMNeT++~6.1 with the INET~4.5 framework. QAD-EDCA is
implemented by subclassing the INET 802.11 HCF/EDCA module chain
(\texttt{Hcf}$\rightarrow$\texttt{Edca}$\rightarrow$\texttt{Edcaf}); the
controller mutates the live EDCAF parameters at runtime.

\textbf{Workload and scenarios.} All schemes share the standards-aligned
traffic model of Table~\ref{tab:traffic}; the scenario matrix is summarized in
Table~\ref{tab:scenarios}. (i) \emph{Steady state}: saturated CBR load at
$N\in\{5,10,15,20\}$, sweeping the high-/low-priority station ratio across four
mixes per $N$ (for $N{=}10$ the VO/VI/BE/BK counts are 1/1/4/4, 2/2/3/3,
3/3/2/2, 4/4/1/1), each 30\,s over 5 seeds; the headline starvation case is
$N{=}10$, 1\,VO/1\,VI/4\,BE/4\,BK. (ii) \emph{Transient}: a 60\,s three-phase
profile calm$[0,20)$/surge$[20,40)$/calm$[40,60)$, where the high-priority load
rises from 1 to 4 VO/VI over a saturated 4\,BE\,$+$\,4\,BK background, 10 seeds.
(iii) \emph{Sensitivity}: the four QAD-EDCA constants
($Q_{th}\!\times\!P_{th}$, $\alpha_{CW}$, $\gamma$, $T_{mon}$) swept at the
headline load. All results are ensemble-averaged. Metrics: per-AC throughput,
delay, and loss; Jain's index; and QoS satisfaction (\AC{VO}/\AC{VI} delay
bounds).

\begin{table}[t]
\centering
\caption{Simulation scenario matrix (shared traffic model, Table~\ref{tab:traffic}).}
\label{tab:scenarios}
\begin{tabular}{@{}llcc@{}}
\toprule
Scenario & Load profile & $N$ & Seeds/time\\
\midrule
Steady state & 4 mixes (VO/VI:BE/BK ratio)    & 5,10,15,20 & 5\,/\,30\,s\\
Transient    & calm\,$\to$\,surge\,$\to$\,calm & 10         & 10\,/\,60\,s\\
Sensitivity  & headline $+$ parameter sweep    & 10         & 5\,/\,30\,s\\
\bottomrule
\end{tabular}
\end{table}

\textbf{Comparison schemes}, all simulated under identical conditions except
where noted:
\begin{enumerate}
\item \textbf{Standard EDCA} --- default parameters (Table~\ref{tab:edca}).
\item \textbf{Tuned Static} --- a hand-tuned static oracle: \AC{BE}
CWmin~$15{\rightarrow}7$/AIFSN~$3{\rightarrow}2$, \AC{BK}
CWmin~$15{\rightarrow}7$/AIFSN~$7{\rightarrow}3$. Represents the best static
configuration \emph{for this load}.
\item \textbf{QAD-EDCA} --- the proposed controller.
\item \textbf{PDCF-DRL} \cite{zuo} --- DRL reference; we cite published numbers
(different setup), as reproducing DRL training is out of scope. Used as an
upper-bound reference, not a head-to-head.
\end{enumerate}

\textbf{Engineering-methodology note.} An early implementation appeared to
``win'' but actually never triggered its own detector. Vector tracing of the
controller state (\texttt{starvationDetected}) revealed it was identically
false: the gain came from a static TXOP misconfiguration
(\AC{BE}/\AC{BK} TXOP set to 2.528~ms instead of the standard 0), which let
\AC{BE}/\AC{BK} burst \emph{and} drained the AP queue so the detector could
never fire. A control experiment (Standard EDCA $+$ the same static TXOP)
reproduced the ``win'' (\AC{BE} 7.22~Mbps $\approx$ the buggy 7.18~Mbps),
proving the effect was static, not adaptive. We corrected the TXOP defaults to
standard (Table~\ref{tab:edca}); afterwards the detector genuinely fires
($\sim$64\% of cycles at the headline load) and the honest \AC{BE} figure is
5.44~Mbps. The lesson---\emph{``dynamic'' must be verified by confirming the
loop closes, not by an improved output number}---shapes the analysis below.

% =====================================================================
\section{Analysis and Discussion}\label{sec:analysis}
% =====================================================================
\subsection{Steady-state starvation relief}
Table~\ref{tab:headline} reports the headline case per AC. QAD-EDCA raises
\AC{BE} throughput $2.15\times$ over Standard (2.53$\rightarrow$5.44~Mbps) and
cuts \AC{BE} loss from 94.7\% to 88.6\%. The real-time \AC{VO}/\AC{VI} flows
(CBR, below capacity) are unaffected---throughput, delay, and loss are
essentially identical across all three schemes. \AC{BK} throughput is likewise
preserved; its mean delay rises modestly under QAD (75$\rightarrow$152~ms),
acceptable for background traffic with no delay bound. This is a clean win over
standard EDCA, achieved with zero manual tuning.

\begin{table}[t]
\centering
\caption{Per-AC steady-state results ($N{=}10$, 1VO/1VI/4BE/4BK).
\AC{VO}/\AC{VI}/\AC{BK} are CBR (below capacity); only \AC{BE} is saturated.}
\label{tab:headline}
\setlength{\tabcolsep}{4pt}
\begin{tabular}{@{}lrrrrrr@{}}
\toprule
 & \multicolumn{3}{c}{Throughput (Mbps)} & \multicolumn{3}{c}{Avg.\ delay (ms)}\\
\cmidrule(lr){2-4}\cmidrule(lr){5-7}
AC & Std & Tuned & QAD & Std & Tuned & QAD\\
\midrule
\AC{VO} & 0.06 & 0.06 & 0.06 & 19.9 & 19.4 & 19.6\\
\AC{VI} & 1.02 & 1.02 & 1.02 & 21.7 & 21.6 & 21.3\\
\AC{BE} & 2.53 & 5.80 & \textbf{5.44} & 1086 & 832 & 957\\
\AC{BK} & 0.16 & 0.16 & 0.16 & 75 & 60 & \textbf{152}\\
\bottomrule
\end{tabular}
\end{table}

\textbf{Scaling.} The advantage grows with network size $N$
(Table~\ref{tab:nscale}, at a fixed 20\% high-priority share): Standard \AC{BE}
throughput collapses from 3.74 to 1.57~Mbps as $N$ goes 5$\to$20, while
QAD-EDCA holds 5.63$\to$4.71~Mbps---widening QAD's gain over Standard from
$1.5\times$ to $3.0\times$.

\begin{table}[t]
\centering
\caption{\AC{BE} throughput (Mbps) vs.\ network size $N$ (20\% high-priority share).}
\label{tab:nscale}
\begin{tabular}{@{}rrrr@{}}
\toprule
$N$ & Standard & Tuned & \textbf{QAD}\\
\midrule
5  & 3.74 & 5.96 & \textbf{5.63}\\
10 & 2.53 & 5.80 & \textbf{5.44}\\
15 & 1.84 & 5.11 & \textbf{4.83}\\
20 & 1.57 & 4.86 & \textbf{4.71}\\
\bottomrule
\end{tabular}
\end{table}

\subsection{The feedback-control ceiling (negative result)}
QAD-EDCA does \emph{not} beat the static optimum. Across all four mixes
(Table~\ref{tab:mixes}) the tuned static configuration exceeds QAD-EDCA by
$\sim$0.3--0.4~Mbps ($\sim$6\%). The gap is stable across every mix and network
size, indicating it is structural rather than a tuning artifact.

\begin{table}[t]
\centering
\caption{\AC{BE} throughput (Mbps) across mixes ($N{=}10$).}
\label{tab:mixes}
\begin{tabular}{@{}lrrr@{}}
\toprule
Mix (VO/VI/BE/BK) & Standard & \textbf{QAD} & Tuned\\
\midrule
1/1/4/4 & 2.53 & \textbf{5.44} & 5.80\\
2/2/3/3 & 2.59 & \textbf{4.82} & 5.20\\
3/3/2/2 & 2.78 & \textbf{4.22} & 4.54\\
4/4/1/1 & 3.32 & \textbf{3.58} & 3.67\\
\bottomrule
\end{tabular}
\end{table}

\emph{Mechanism.} The ceiling follows directly from the observability trap of
Sec.~\ref{sec:problem}. When QAD-EDCA lowers \AC{BE} CWmin it succeeds in
draining the AP queue; the occupancy predicate \eqref{eq:starv} then reads
``not starving,'' so the boost decays under recovery. Empirically the trigger
rate falls from 64\% (relief OFF in CWmin) to $\sim$15\% once bidirectional
CWmin is enabled, and hysteresis cannot hold it because the AP downlink queue
is bursty and instantaneous samples are mostly empty. To beat a feedforward
optimum in steady state a feedback controller would have to hold maximum boost
\emph{unconditionally}---at which point it has become a static scheme and
abandoned adaptivity. This is the central honest insight of the work: in
saturated steady state, feedback cannot exceed good feedforward.

\subsection{Robustness to algorithm constants}
A four-way sweep of $\alpha_{CW}$, $\gamma$, $T_{mon}$, and the
($Q_{th},P_{th}$) pair changes steady-state \AC{BE} throughput by
$\le1.5\%$ (Table~\ref{tab:sens}). Under sustained saturation the controller
drives parameters to their operating bounds regardless of the constants, so any
reasonable setting converges to the same effect---i.e.\ no per-deployment
tuning of QAD-EDCA's own constants is needed. (The flip side: the constants
matter only in the trigger and transient phases, which the dynamic experiment
below exercises.)

\begin{table}[t]
\centering
\caption{Sensitivity of steady-state \AC{BE} throughput to QAD-EDCA's own
constants ($N{=}10$, headline load).}
\label{tab:sens}
\begin{tabular}{@{}llr@{}}
\toprule
Constant & Swept range & \AC{BE} spread (Mbps)\\
\midrule
$\alpha_{CW}$    & 1.5--4.0             & 5.43--5.44 (0.2\%)\\
$\gamma$         & 0.1--0.7             & 5.41--5.49 (1.5\%)\\
$T_{mon}$        & 50--500\,ms          & 5.41--5.45 (0.8\%)\\
$Q_{th},P_{th}$  & 0.5--0.9 / 0.05--0.2 & 5.41--5.45 (0.7\%)\\
\bottomrule
\end{tabular}
\end{table}

\subsection{Dynamic load: the value of adaptivity}
To probe the regime where adaptivity matters we use a 60~s,
three-phase load: calm~$[0,20)$, surge~$[20,40)$ (high-priority load rises
from 1 to 4 VO/VI), calm~$[40,60)$, over a saturated 4\,BE\,$+$\,4\,BK
background; results are ensemble-averaged over 10 seeds.
Table~\ref{tab:dynamic} gives per-phase means.

\begin{table}[t]
\centering
\caption{Dynamic load, per-phase ensemble means.}
\label{tab:dynamic}
\begin{tabular}{@{}llrrr@{}}
\toprule
Scheme & Phase & \AC{BE} (Mbps) & $d_{VO}$ (ms) & $d_{VI}$ (ms)\\
\midrule
Standard & calm  & 2.6 & 1.1  & 2.6\\
Standard & surge & 1.9 & 28.7 & 33.1\\
Tuned    & calm  & 6.0 & 1.4  & 3.1\\
Tuned    & surge & 4.0 & 2.3  & 4.5\\
\textbf{QAD} & calm  & \textbf{5.6} & 1.3 & 2.8\\
\textbf{QAD} & surge & \textbf{3.7} & 4.0 & 7.1\\
\bottomrule
\end{tabular}
\end{table}

Two honest readings follow. (i) \emph{QAD-EDCA tracks the load without manual
tuning}, staying within $\sim$6\% of the Tuned oracle throughout while needing
no foreknowledge---when the surge arrives it adapts, and it relaxes afterward.
(ii) \emph{Standard EDCA fails the real-time flows under surge:} although its
\emph{phase-mean} \AC{VO}/\AC{VI} delay (28.7/33.1~ms) looks bounded, the
time series spikes to $\approx$240~ms during the surge---exceeding the 150~ms
\AC{VO} budget---whereas Tuned and QAD-EDCA keep instantaneous delay low
(Fig.~\ref{fig:hidelay}). QAD-EDCA thus recovers most of the oracle's benefit
on \emph{both} axes with no configuration. It does not, however, beat Tuned on
any axis here, consistent with the steady-state ceiling.

\begin{figure}[t]
\centering
\includegraphics[width=\columnwidth]{fig_dynamic_hi_delay.pdf}
\caption{High-priority delay under the transient surge. Standard EDCA spikes
to $\approx$240~ms (above the 150~ms \AC{VO} budget); Tuned and QAD-EDCA hold
it low. QAD-EDCA achieves this with no manual tuning.}
\label{fig:hidelay}
\end{figure}

\begin{figure}[t]
\centering
\includegraphics[width=\columnwidth]{fig_dynamic_be_throughput.pdf}
\caption{\AC{BE} throughput over time. QAD-EDCA tracks the Tuned oracle
($\sim$2.15$\times$ Standard) across calm and surge without prior tuning.}
\label{fig:bethr}
\end{figure}

\subsection{Honest limitations}
The detector observes the \emph{AP downlink} queue. In the
station\,$\rightarrow$\,AP\,$\rightarrow$\,sink topology, starvation often
forms on the \emph{station uplink}, so AP-side occupancy fires only when enough
saturated \AC{BE}/\AC{BK} stations make the AP downlink the bottleneck (e.g.\
1VO/1VI/4BE/4BK triggers; 3VO/3VI/2BE/2BK does not). This is an intrinsic
limit of AP-side measurement, not a defect, and it is the main reason
feedback's reach is bounded here. An uplink/airtime-based indicator is the
natural next step and may also relax the steady-state ceiling.

% =====================================================================
\section{Comparison with State-of-the-Art}\label{sec:compare}
% =====================================================================
\begin{figure}[t]
\centering
\includegraphics[width=\columnwidth]{positioning_tradeoff.pdf}
\caption{Performance--complexity positioning. QAD-EDCA occupies the
no-training, no-foreknowledge quadrant between static EDCA and DRL.}
\label{fig:pos}
\end{figure}

Table~\ref{tab:sota} places QAD-EDCA among representative schemes
(Standard, AC-agnostic DRL, the static oracles, and AC-aware DRL).

\begin{table}[t]
\centering
\caption{Positioning. ``Foreknowledge'' = needs the load known in advance;
``Training'' = needs offline learning.}
\label{tab:sota}
\setlength{\tabcolsep}{3.5pt}
\begin{tabular}{@{}lcccc@{}}
\toprule
 & Std. & \textbf{QAD} & Static & PDCF\\
 & EDCA & \textbf{EDCA} & oracle & DRL\\
\midrule
Adaptive at runtime & no & \textbf{yes} & no & yes\\
Needs foreknowledge & no & \textbf{no} & \emph{yes} & no\\
Needs training      & no & \textbf{no} & no & \emph{yes}\\
Per-cycle cost      & $O(1)$ & \textbf{$O(1)$} & $O(1)$ & $O(\text{DRL})$\\
\AC{BE} relief      & none & \textbf{2.15$\times$} & $\sim$2.3$\times$ & near-equal\\
\bottomrule
\end{tabular}
\end{table}

\textbf{Among plug-and-play methods}---those that need neither foreknowledge of
the load nor training---QAD-EDCA is the best: it doubles starved \AC{BE}
throughput with no QoS loss and acts from the first monitoring cycle at $O(1)$
cost. The only schemes that beat QAD-EDCA are the two ``need-to-know-first''
upper bounds---the static oracle (needs the load) and DRL (needs training)---which
is exactly the gap QAD-EDCA is designed to fill, and the performance--complexity
trade-off is visualized in Fig.~\ref{fig:pos}.

% =====================================================================
\section{Conclusion}
% =====================================================================
We presented QAD-EDCA, a lightweight $O(1)$ queue-aware controller that
relieves IEEE~802.11 EDCA starvation, raising starved \AC{BE} throughput
$2.15\times$ over standard EDCA with no QoS loss and zero tuning or training.
We also reported, honestly, what it cannot do: in saturated steady state a
queue-feedback controller cannot surpass a well-chosen feedforward static
optimum, because relieving starvation removes its own trigger---a result we
confirmed against a well-tuned static optimum and a mechanistic explanation,
and which we reached only after correcting a ``fake dynamics'' confound exposed
by vector tracing. QAD-EDCA is therefore best understood as the strongest
\emph{plug-and-play} option---no foreknowledge, no training---trailing only
foreknowledge- and training-dependent upper bounds. Future work: an
uplink/airtime indicator (to extend reach
and possibly relax the ceiling), genuinely non-stationary and heterogeneous
loads, and using QAD-EDCA as a warm start during DRL convergence.

% =====================================================================
\begin{thebibliography}{99}
\bibitem{ieee80211} IEEE Std 802.11-2020, ``IEEE Standard for Information
Technology---Wireless LAN MAC and PHY Specifications,'' IEEE, 2020.
\bibitem{ugwu} G.~O. Ugwu, U.~N. Nwawelu, and M.~A. Ahaneku, ``Effect of
service differentiation on QoS in IEEE 802.11e EDCA: a simulation approach,''
\emph{J. Eng. Appl. Sci.}, vol.~69, no.~1, pp.~1--15, 2022.
\bibitem{mammeri} S.~Mammeri, M.~Yazid, R.~Kacimi, and L.~Bouallouche-Medjkoune,
``Starvation avoidance-based dynamic multichannel access for low priority
traffics in 802.11ac,'' \emph{Comput. Electr. Eng.}, vol.~90, art.~106942, 2021.
\bibitem{tuan} Y.~P. Tuan \emph{et al.}, ``Improving QoS mechanisms for IEEE
802.11ax with overlapping basic service sets,'' \emph{Wireless Netw.},
vol.~29, pp.~387--401, 2023.
\bibitem{yi} P.~Yi, W.~Cheng, J.~Wang, J.~Pan, Y.~Ouyang, and W.~Zhang,
``Intelligent Multi-link EDCA Optimization for Delay-Bounded QoS in Wi-Fi~7,''
arXiv:2509.25855, 2025.
\bibitem{zuo} Z.~Zuo, D.~Wang, X.~Nie, X.~Pan, M.~Deng, and M.~Ma, ``PDCF-DRL:
a contention window backoff scheme based on deep reinforcement learning for
differentiating access categories,'' \emph{J. Supercomput.}, vol.~81,
art.~213, 2025.
\bibitem{du} X.~Du, X.~Fang, R.~He, L.~Yan, L.~Lu, and C.~Luo, ``Federated deep
reinforcement learning-based intelligent channel access in dense Wi-Fi
deployments,'' arXiv:2409.01004, 2024.
\bibitem{li} Q.~Li, B.~Lv, Y.~Hong, and R.~Wang, ``ReinWiFi: Application-layer
QoS optimization of WiFi networks with reinforcement learning,''
arXiv:2405.03526, 2024.
\bibitem{jain} R.~Jain, D.~Chiu, and W.~Hawe, ``A quantitative measure of
fairness and discrimination for resource allocation in shared computer
systems,'' DEC Research Report TR-301, 1984.
\bibitem{ax} IEEE Std 802.11ax-2021, ``Amendment 1: Enhancements for
High-Efficiency WLAN,'' IEEE, 2021.
\end{thebibliography}

\end{document}
